В начале необходимо определить детали реализации, такие как размер рабочей сетки и основные параметры, которые будут использоваться в вычислениях.


Пусть у нас будет простая двумерная квадратная сетка размером $N \times N$. Длина стороны каждой ячейки равна $L$. Следуя классическому варианту алгоритма каждая ячейка содержит определенные в ее центре вектора $u_0(x,y)$ и $u_1(x,y)$ \cite{stamforgames}. В каждый момент времени мы будем вычислять новые значения $u_1(x,y)$ на основе предыдущего значения $u_0(x,y)$, используя четыре этапа алгоритма Stable Fluids.

То есть мы можем представить векторное поле как два двумерных массива размером $N \times N$, где каждый элемент массива содержит вектор в соответствующей ячейке сетки.

Таким образом, мы представляем нашу симуляцию как набор буферов, хранящих данные о скорости и других физических величинах для каждой ячейки сетки. Эти буферы будут использоваться в вычислительных шейдерах для выполнения каждого из этапов алгоритма.

\subsection{Инициализация симуляции}

\subsubsection{Параметры симуляции}

Определим основные константы:

\begin{itemize}
    \item Размерность: $N_{DIM}=2$. 2D симуляция.
    \item Начало координат: $O=(0,0)$. Центр области отрисовки.
    \item Физический размер пространства: $L=(750,750)$ в пикселях. Область отрисовки
    \item Число ячеек: $N_x=N_y=N=250$.
    \item Размер ячейки: $\Delta x=\Delta y=L_x/N_x=L_y/N_y=3$.
    \item Шаг по времени: $\Delta t=1/60\ \text{s}$. 60 кадров в секунду.
    \item Итерации решения давления: $\text{PRESSURE\_ITERATIONS}=20$.
    \item Итерации решения вязкости: $\text{VISCOSITY\_ITERATIONS}=10$.
    \item Размер группы (для шейдера): $\text{WORKGROUP\_SIZE}=16$. 16x16=256 потоков.
\end{itemize}

\subsubsection{Динамические параметры}

Для обеспечения интерактивности симуляции, мы будем использовать динамические параметры, которые могут быть изменены пользователем во время выполнения симуляции.

Определим начальные значения динамических параметров симуляции:

\begin{itemize}
    \item Вязкость: $\nu=\text{VISCOSITY}=0.005$.
    \item Рассеивание: $\text{DISSIPATION}=0.999$.
    \item Коэффициент завихренности: $\text{VORTEX\_STRENGTH}=25$.
    \item Гравитация (ось $y$): $\text{GRAVITY}=-20$.
    \item Количество добавляемого вещества: $\text{FAUCET\_DYE}=200$.
    \item Позиция мыши: $\text{MOUSE\_POS}=(N_x/2,N_y/2)$.
    \item Скорость мыши: $\text{MOUSE\_VEL}=(0,0)$.
\end{itemize}

Для передачи этих параметров на GPU, используются uniform буферы, которые позволят нам передавать эти значения в шейдеры на каждом кадре симуляции.

\subsubsection{Буферы}

Для удобства реализации на GPU, мы будем использовать буферы для хранения данных сетки, такие как буфер скорости, буфер давления и буфер вещества-краски (dye). Каждый из этих буферов будет представлен в виде пар массивов, который может быть эффективно обработан параллельно на GPU при помощи техник ping-pong (переключение между двумя буферами для чтения и записи). Также для решения уравнения Пуассона потребуется буфер для дивергенции.



\subsection{Создание шейдеров}

Для реализации каждого этапа алгоритма Stable Fluids \eqref{eq:stages_pipeline}, были созданы отдельные вычислительные шейдеры на языке WGSL.

Ниже приведен полный список шейдеров, используемых в реализации симуляции жидкости:

\begin{description}
    \item[00\_header.wgsl] Общие определения, константы и функции, подключаемые в другие шейдеры.
    \item[01\_render.wgsl] Шейдер вывода: преобразует данные симуляции в финальный цвет для экрана.
    \item[11\_add\_force.wgsl] Добавляет внешние силы в поле скоростей (input velocity += force).
    \item[12\_vortex\_force.wgsl] Применяет вихревые силы для сохранения завихрений в потоке.
    \item[13\_gravity.wgsl] Моделирует действие гравитации на плотность/скорость среды.
    \item[21\_advect\_vel.wgsl] Адвекция скорости — перенос значений скорости по самому полю.
    \item[31\_viscosity.wgsl] Учет вязкости: сглаживание и диффузия скоростей.
    \item[41\_clear\_pressure.wgsl] Сбрасывает или инициализирует поле давления перед итерациями решения.
    \item[42\_divergence.wgsl] Вычисляет дивергенцию скоростного поля для контроля несжимаемости.
    \item[43\_pressure\_jacobi.wgsl] Итеративный решатель давления методом Якоби для устранения дивергенции.
    \item[44\_subtract\_gradient.wgsl] Корректирует скорости вычитанием градиента давления (проекция на несжимаемое поле).
    \item[51\_add\_dye.wgsl] Добавляет краску/плотность в поле (внешние источники плотности).
    \item[52\_splat\_dye.wgsl] Локальные всплески краски (splat) — быстрые импульсные добавления плотности.
    \item[53\_advect\_dye.wgsl] Адвекция краски по скоростному полю (перенос плотности).
\end{description}

Заметим, что первая цифра в названии шейдера указывает на этап алгоритма, а вторая \textemdash{} на конкретную операцию внутри этого этапа. Такой подход упрощает организацию и понимание структуры шейдеров.

Этап 0 \textemdash{} рендеринг, этапы 1-4 \textemdash{} основные этапы алгоритма Stable Fluids, этап 5 \textemdash{} работа с краской (dye).

% ...existing code...